\subsection{Motivación}

La resolución de los problemas descritos en la sección anterior no es un
ejercicio puramente teórico. Las limitaciones del tool calling convencional
sobre MCP tienen consecuencias tangibles que justifican una investigación
sistemática desde múltiples perspectivas.

\subsubsection{Impacto económico}

En aplicaciones de producción, cada token procesado por un LLM tiene un costo
monetario directo. Los proveedores de modelos ---OpenAI, Anthropic, Google,
Amazon--- facturan por millón de tokens de entrada y de salida, con precios que
varían según el modelo pero que en todos los casos escalan linealmente con el
consumo. Cuando un agente conectado a múltiples servidores MCP carga
estáticamente 150,000 tokens de definiciones de herramientas en cada
interacción, ese costo se repite en \textit{cada solicitud del usuario},
independientemente de cuántas herramientas se utilicen realmente. A escala
---miles de usuarios, millones de solicitudes diarias--- la diferencia entre
una arquitectura que consume 150,000 tokens y una que consume 2,000 tokens por
interacción puede representar órdenes de magnitud en costos operativos. Sin
embargo, no existe una cuantificación rigurosa de estos ahorros a través de
diferentes tipos de tareas y configuraciones, lo cual impide que las
organizaciones tomen decisiones informadas sobre su arquitectura.

\subsubsection{Ausencia de guías para desarrolladores}

El ecosistema de desarrollo de sistemas agénticos carece de criterios empíricos
para elegir entre llamadas directas y ejecución de código. Un desarrollador que
diseña un sistema basado en MCP enfrenta una decisión arquitectónica
fundamental al inicio del proyecto: ¿debe su agente llamar herramientas
directamente, o debe generar código que las orqueste? Esta decisión tiene
implicaciones profundas en el diseño del sistema, la infraestructura necesaria
(sandbox de ejecución, seguridad, monitoreo) y el comportamiento en producción.
Actualmente, esta elección se basa en intuición, anécdotas o en las
recomendaciones de un único artículo técnico \citep{anthropic2024codeexec}, sin
un marco cuantitativo que relacione las características de la tarea con el modo
de ejecución óptimo.

\subsubsection{Crecimiento acelerado del ecosistema MCP}

Desde su introducción en 2024, el Model Context Protocol ha experimentado una
adopción acelerada. Empresas como Anthropic, Amazon Web Services, Microsoft y
múltiples startups han construido servidores MCP para sus plataformas, y la
comunidad open-source ha contribuido cientos de implementaciones que cubren
desde bases de datos hasta APIs de servicios en la nube
\citep{anthropic2024mcp}. Esta proliferación de servidores y herramientas
agrava directamente el problema de la carga estática: conforme un agente se
conecta a más servidores, el volumen de definiciones de herramientas en su
contexto crece proporcionalmente. Lo que hoy es un problema manejable con un
puñado de herramientas se convertirá en un cuello de botella crítico conforme
el ecosistema madure y los agentes necesiten acceso a decenas de servidores
simultáneamente.

\subsubsection{Vacío en la literatura académica}

A pesar de la rápida adopción industrial del protocolo MCP y de los patrones de
ejecución de código, la literatura académica no ha abordado de manera
sistemática la comparación entre estos dos paradigmas de ejecución. Existen
benchmarks para evaluar la capacidad general de tool-use de los LLMs, así como
evaluaciones de sistemas multi-agente en tareas específicas, pero ninguno se
enfoca en el eje particular de \textit{cómo} se ejecutan las herramientas y su
impacto en eficiencia. Esta ausencia de evidencia empírica deja un vacío que
este trabajo busca llenar mediante un diseño experimental controlado,
reproducible y extensible por la comunidad.
